%==============================================================================
% ACKNOWLEDGEMENTS
%==============================================================================

\chapter*{Acknowledgements}
\addcontentsline{toc}{chapter}{Acknowledgements}

This work was carried out under conditions that would have been impossible
without the benevolence of one Julian Oskar Matschinske. I also want to extend
my gratitude to Peter Haferl, Irakli Lezhava, Mihail Boyadzhiev, Alexander
Held, Izuka Kai, and Michael Lauber for the discussions that propelled this
work. The cost of food, shelter, and the time required to think has been borne
by friends who chose, repeatedly and without expectation of return, to keep me
going. They have my gratitude in a form that exceeds what I can write.

I want to single out one debt specifically. My dog, Merovingian Visconti
Rembrandt Moriarty Picollo, for being the most supportive individual and
constant companion of this work. Through every chapter, every derivation, every
long evening when the framework was either coming together or refusing to, the
dog has been there --- uninterested in the outcome, unbothered by the
difficulty, requiring only the basic things that dogs require and offering, in
return, the steady presence that makes sustained thinking possible. The
framework was developed with the dog at my feet or nearby. Whatever this
monograph contains was thought through in that company.

It is not customary to acknowledge a dog in a scientific monograph. I am
acknowledging mine because the customs in this matter are wrong. A mind that
wanders far enough to produce work like this needs an anchor that does not
wander, and dogs are unusually good at being that anchor. The work would not
exist without the conditions that allowed it, and the dog was one of those
conditions.

\bigskip
\noindent
To my friends: thank you for buying me the time.\\
To my dog: thank you for keeping me here while I used it.
